% GENERATED by code/python/weekly_report.py from the template templates/DATA.tex for the
% sweep order 132; do not edit by hand: edit templates/DATA.tex and generate again.
\documentclass[11pt]{article}
\usepackage[margin=1in]{geometry}
\usepackage{amsmath,amssymb,amsfonts}
\usepackage{bm}
\usepackage{amsthm,mathrsfs}
\newtheorem{lemma}{Lemma}[section]
\theoremstyle{remark}
\newtheorem{remark}[lemma]{Remark}
\usepackage{enumitem}
\usepackage{xurl}
\usepackage[colorlinks=true,linkcolor=blue,citecolor=blue,urlcolor=blue]{hyperref}

% file and column names: typewriter, breakable at any character (xurl)
\newcommand{\src}[1]{\path{#1}}
% a file of Section 1, by its label: F<number>
\newcommand{\fref}[1]{\ref{#1}}

% TEMPLATE: code/python/weekly_report.py writes notes/DATA_<order>.tex from this file (notes/templates/DATA.tex), with
% run_values replaced by run_values_<order> (<order> the sweep order of the run); compile that
% one. Edit the text here.
% the values of one run of the weekly simulation, generated by code/python/weekly_report.py
% ============================================================================================
% THIS FILE IS GENERATED AUTOMATICALLY. DO NOT EDIT IT BY HAND: EVERY CHANGE IS LOST
% THE NEXT TIME IT IS GENERATED.
%
% IT IS WRITTEN BY code/python/weekly_report.py (RUN BY weekly_simulation/run.sh, OR BY
%     python3 code/python/weekly_report.py weekly_simulation/results/<tag> [nofigures]
% ) FROM THE RESULTS OF ONE RUN WITH SWEEP ORDER 132, AND OVERWRITTEN BY THE NEXT ONE.
% THIS VERSION: weekly_simulation/results/Soln_labeled__Vax2021__2021__132/, 2026-11-22 19:26.
%
% IT IS INCLUDED BY notes/DATA_132.tex, notes/simulation_weekly_optimal_132.tex and notes/stratified_seir_ngm.tex
% (IF ITS PREAMBLE INCLUDES run_values_132). RERUN pdflatex ON THEM AFTER IT IS GENERATED AGAIN.
%
% TO CHANGE A VALUE OR ITS WORDING, CHANGE weekly_report.py (FUNCTIONS values AND
% delta_values) AND GENERATE THE FILE AGAIN. A NEW \Run... MACRO USED IN THE NOTES MUST
% BE ADDED THERE TOO, OR THE NOTES DO NOT COMPILE.
% ============================================================================================
\newcommand{\RunTag}{Soln\_\allowbreak{}labeled\_\allowbreak{}\_\allowbreak{}Vax2021\_\allowbreak{}\_\allowbreak{}2021\_\allowbreak{}\_\allowbreak{}132}
\newcommand{\RunInitialFile}{data/\allowbreak{}real/\allowbreak{}Soln\_\allowbreak{}labeled.csv}
\newcommand{\RunDosesFile}{data/\allowbreak{}real/\allowbreak{}Vax2021.xlsx}
\newcommand{\RunParams}{2021}
\newcommand{\RunTzero}{347}
\newcommand{\RunTzeroDate}{12 December 2020}
\newcommand{\RunOrder}{$\chi_1$, $\chi_3$, $\chi_2$}
\newcommand{\RunTvacc}{354}
\newcommand{\RunTvaccDate}{19 December 2020}
\newcommand{\RunWeeks}{55}
\newcommand{\RunTlast}{732}
\newcommand{\RunTend}{739}
\newcommand{\RunTendDate}{8 January 2022}
\newcommand{\RunDaysNoVacc}{7}
\newcommand{\RunFigDir}{../weekly_simulation/results/Soln_labeled__Vax2021__2021__132/figures/}
\newcommand{\RunSwitch}{$t=398$, $t=460$ and $t=548$ (1 February 2021, 4 April 2021, 1 July 2021)}
\newcommand{\RunRhoOptAbove}{21}
\newcommand{\RunCapWeeks}{44}
\newcommand{\RunVPeak}{94.4}
\newcommand{\RunVPeakDate}{29 May 2021}
\newcommand{\RunInfPeak}{20.7}
\newcommand{\RunInfPeakDate}{23 October 2021}
\newcommand{\RunHospPeak}{1.29}
\newcommand{\RunHospPeakDate}{30 October 2021}
\newcommand{\RunDeathsStart}{132}
\newcommand{\RunDeathsEnd}{921}
\newcommand{\RunSeroPct}{16}
\newcommand{\RunDosesOne}{223}
\newcommand{\RunDosesTwo}{182}
\newcommand{\RunDosesThree}{69}
\newcommand{\RunInitialText}{the initial state at $t=\RunTzero$, \texttt{\RunInitialFile}; $\RunSeroPct\%$ seropositive}
\newcommand{\RunDosesText}{the weekly doses by location and dose, \texttt{\RunDosesFile}: \RunWeeks{} weeks, with $\RunDosesOne$, $\RunDosesTwo$ and $\RunDosesThree$ million first, second and third doses}
\newcommand{\RunOtherInputs}{not from public sources}
\newcommand{\CmpEigenvectorDeathsEnd}{750}
\newcommand{\CmpEigenvectorRhoAbove}{31}
\newcommand{\CmpEigenvectorInfPeak}{15.5}
\newcommand{\CmpEigenvectorHospPeak}{1.03}
\newcommand{\RunOmegaInv}{182.5}
\newcommand{\RunRMin}{0.714}
\newcommand{\RunRMax}{0.985}
\newcommand{\RunEpsFirst}{0.9}
\newcommand{\RunDeltaBoundLo}{7.9\cdot10^{-4}}
\newcommand{\RunDeltaBoundHi}{3.4\cdot10^{-3}}
\newcommand{\RunEpsFirstTo}{ in the weeks before $t=403$}
\newcommand{\RunDeltaLast}{ (with $\varepsilon_1=0.7$, in the weeks from $t=550$: between $4.3\cdot10^{-3}$ and $9.6\cdot10^{-3}$)}
\newcommand{\RunDeltaText}{in 12 of the 55 weeks, the first at $t=361$}
\newcommand{\CmpOptimalDeathsEnd}{921}
\newcommand{\CmpOptimalRhoAbove}{21}
\newcommand{\CmpOptimalInfPeak}{20.7}
\newcommand{\CmpOptimalHospPeak}{1.29}


\title{The age- and location-stratified SEIR model: the data\thanks{Parts of this
manuscript were edited with the assistance of an AI tool (Claude Opus 5.5, Anthropic).
The authors are solely responsible for all proofs and statements.}}
\author{Zhilan \and John \and Yuliya \and Ludmil}
\date{}

\begin{document}
\maketitle

\begin{abstract}
The data of the model: the public sources, the scripts that build the inputs from them,
and other data as described below.
\end{abstract}

\section{Data}\label{sec-data}
There are $m=51$ locations (the 50 states and D.C., in the order of their FIPS codes)
and $n=16$ age groups ($0$--$4$, \dots, $70$--$74$, $75+$).

\subsection{Files}\label{sec-files}

\subsubsection{Downloaded}\label{sec-downloaded}
Kept as downloaded, in
\src{MathBio/data/Downloaded_from_source/}; paths relative to this directory.
\begin{enumerate}[label=F1-\arabic*, ref=F1-\arabic*, series=files, leftmargin=*]
\item\label{f-jtb} \src{papers/2026FengDayGlasser.pdf}: the model paper~\cite{2026Feng_Day_Glasser}.
\item\label{f-jtbs} \src{papers/2026FengDayGlasserSupplement.pdf}: its supplement (Tables S1--S4).
Sections~\ref{sec-age}, \ref{sec-const}, \ref{sec-var}.
\item\label{f-g23} \src{papers/2023GlasserFengVoJonesClarke.pdf}:~\cite{2023Glasser_Feng}; its
Fig.~4 is digitized for $\beta^0$. Sections~\ref{sec-time}--\ref{sec-inputs}.
\item\label{f-f17} \src{papers/2017FengHillCurnsGlasser.pdf}:~\cite{2017Feng_Hill_Curns_Glasser}.
Section~\ref{sec-loc}.
\item\label{f-prem} \src{papers/2017PremCookJit.pdf}:~\cite{2017Prem_Cook_Jit}, and its
supporting text S1 \src{papers/2017PremCookJitSupplementS1.pdf}. Section~\ref{sec-contacts}.
\item\label{f-baj}
\src{papers/2021BajemaWiegandCuffePatelIachanLimLeeMoyseHaversHardingFryHallMartinBielDengMeyerMathurKyleGundlapalliThornburgPetersenEdens.pdf}:~\cite{2021Bajema_Wiegand},
and its online supplement (eMethods; eTables 3--6, the seroprevalence by
jurisdiction and age group in the four periods)
\src{papers/2021BajemaWiegandCuffePatelIachanLimLeeMoyseHaversHardingFryHallMartinBielDengMeyerMathurKyleGundlapalliThornburgPetersenEdensSupplement.pdf}.
Section~\ref{sec-inputs}.
\item\label{f-kingman} \src{papers/1961Kingman.pdf}:~\cite{1961Kingman-aa}, the
log-convexity of the spectral radius used in \src{stratified_seir_ngm.tex}
(Lemma~5.1); not a data source.
\item\label{f-sc} \src{census/sc-est2024-alldata6.csv.gz}: state population estimates by age,
Vintage 2024~\cite{2024Census_scest2024}. Section~\ref{sec-pop}.
\item\label{f-co} \src{census/co-est2024-alldata.csv.gz}: county population estimates,
Vintage 2024~\cite{2024Census_coest2024}. Section~\ref{sec-loc}.
\item\label{f-gaz} \src{census/2024_Gaz_counties_national.zip}: county land areas,
2024 Gazetteer~\cite{2024Census_gazetteer}. Section~\ref{sec-loc}.
\item\label{f-cenpop} \src{census/CenPop2010_Mean_ST.txt.gz}: 2010 centers of population of
the states~\cite{2010Census_cenpop}. Section~\ref{sec-loc}.
\item\label{f-premx} \src{prem2017/pcbi.1005697.s002.zip}: supporting file S2
of~\cite{2017Prem_Cook_Jit}, the contact matrices; the one used is
\src{MUestimates_all_locations_2.xlsx} in it. Section~\ref{sec-contacts}.
\item\label{f-ox} \src{oxcgrt/OxCGRT_compact_subnational_v1.csv.gz}: the final OxCGRT data
set~\cite{2021Hale_OxCGRT,2023OxCGRT_data}. Section~\ref{sec-loc}.
\item\label{f-sero} \src{cdc/seroprevalence.csv.gz}: the CDC Nationwide Commercial Laboratory
Seroprevalence Survey~\cite{2023CDC_sero}. Section~\ref{sec-inputs}.
\item\label{f-var} \src{cdc/variants_usa.csv.gz}: CDC SARS-CoV-2 Variant
Proportions, USA, weighted estimates~\cite{2023CDC_variants}. Section~\ref{sec-var}.
\end{enumerate}
Text files are gzipped (\src{.gz}) and zip archives kept as they are; the scripts read
both without unpacking.

\subsubsection{Built}\label{sec-built}
The inputs of the model, in \src{MathBio/code/python/data/},
each written by the script named, from the files listed.
\begin{enumerate}[label=F1-\arabic*, ref=F1-\arabic*, resume=files, leftmargin=*]
\item\label{f-N} \src{census/N_51x16.csv}: $N_{ij}$; \src{census_N_51x16.py}
from~\fref{f-sc}. Section~\ref{sec-pop}.
\item\label{f-U} \src{prem2017/U_16x16.csv}: $U$; \src{prem_U_16x16.py}
from~\fref{f-premx}. Section~\ref{sec-contacts}.
\item\label{f-d} \src{census/d_51x51.csv}: $d_{ip}$; \src{geo_51.py}
from~\fref{f-cenpop}. Section~\ref{sec-loc}.
\item\label{f-L} \src{census/L_51.csv}: $\rho_i$ and $L_i$; \src{geo_51.py}
from~\fref{f-co} and~\fref{f-gaz}. Section~\ref{sec-loc}.
\item\label{f-C} \src{model/C_816x816.csv.gz}: the contact matrix $C$ (gzipped CSV; rows
$(i,j)$ and columns $(p,q)$ labelled by state and age group); \src{build_C_51x16.py}
from~\fref{f-U}, \fref{f-d}, \fref{f-L}. Section~\ref{sec-loc}.
\item\label{f-s} \src{oxcgrt/s_51_daily.csv}: $s_i(t)$, daily; \src{oxcgrt_s_51.py}
from~\fref{f-ox}. Section~\ref{sec-loc}.
\item\label{f-t2} \src{glasser2023/table2.csv}: $p_j$, $d_j$, $\mathrm{IFR}_j$, $c_j^0$, typed
from Table~2 of~\fref{f-g23}; \src{glasser2023_inputs.py}. Section~\ref{sec-age}.
\item\label{f-beta} \src{glasser2023/beta0.csv}: $\beta_j^0$, digitized from Fig.~4
of~\fref{f-g23}; \src{glasser2023_inputs.py}. Section~\ref{sec-age}.
\item\label{f-vd} \src{cdc/variant_days.csv}: the days of the variants;
\src{cdc_variants.py} from~\fref{f-var}. Section~\ref{sec-var}.
\item\label{f-x0} \texttt{\RunInitialFile}: the initial state of the computation of
Section~\ref{sec-comp}; the initial conditions. CSV,
$1+13m=664$ rows and $2+n=18$ columns. Row 1: ``variable'', ``location'', then the age
groups 0--4, \dots, 75+. Then 13 blocks of $m$ rows, the compartments in the order $S$,
$P$, $E_s$, $E_p$, $V$, $I_a$, $I_{ps}$, $I_{pm}$, $I_s$, $I_m$, $I_h$, $R$, $M$; within a
block the locations in FIPS order; column A the compartment, column B the location, then
the $n$ age groups. The program checks the labels of row 1 and of columns A and B
(\src{read_state} in \src{seir7_ode.py}). Section~\ref{sec-inputs}.
\item\label{f-doses} \texttt{\RunDosesFile}: the weekly doses of the computation of
Section~\ref{sec-comp}. 
\end{enumerate}

\subsection{Time}\label{sec-time}

The time $t$ is the number of days from 31 December 2019. The model is started at the day
$t_0$ of the initial state, a parameter of the set: $t_0=251$, 7 September
2020~\cite[Section~3]{2023Glasser_Feng}, in the set of~\cite{2026Feng_Day_Glasser}, and
$t_0=347$, 12 December 2020, in the set of Unknown 2020/21; the set \texttt{\RunParams} starts it
at $t_0=\RunTzero$, \RunTzeroDate.
Vaccination began on 14 December 2020~\cite[Section~2.2]{2026Feng_Day_Glasser},
$t=349$; the parameter set \texttt{\RunParams}, used in the computation of
Section~\ref{sec-comp}, starts it at $t_{\mathrm{vacc}}=\RunTvacc$, \RunTvaccDate{}
(Section~\ref{sec-params-const}).

\subsection{Population}\label{sec-pop}

$N_{ij}$ is the resident population of location $i$ and age group $j$ on July 1,
2020, of the Vintage 2024 estimates~\cite{2024Census_scest2024}
(\fref{f-sc}): column \src{POPESTIMATE2020}, both sexes, all races and origins,
summed over single years of age (85+ into 75+). Total: $331\,577\,720$.
(\fref{f-N})

\subsection{Contacts}\label{sec-contacts}

$U$ is the projected contact matrix of~\cite{2017Prem_Cook_Jit} for the United
States, all locations (\fref{f-premx}, \src{MUestimates_all_locations_2.xlsx},
sheet ``United States of America''), balanced as
in~\cite[Section~2.4]{2023Glasser_Feng} by $U_{jq}=\sqrt{M_{jq}M_{qj}}$. Row sums:
$2.2$--$22.1$ contacts per day. (\fref{f-U})

\subsection{Locations}\label{sec-loc}

\subsubsection{Distances}\label{sec-dist}
$d_{ip}$ is the distance in miles between the 2010
centers of population~\cite{2010Census_cenpop} (\fref{f-cenpop}) of locations $i$
and $p$, as in~\cite[Fig.~6]{2023Glasser_Feng}: the geodesic distance on the
WGS84 ellipsoid~\cite{1975Vincenty}. (\fref{f-d})

\subsubsection{Activity factors}\label{sec-L}
$L_i=\ln\rho_i/\ln\rho$~\cite[Section~2.4]{2023Glasser_Feng},
densities in persons per square mile of land. As
in~\cite[Supplement, Section~2a]{2026Feng_Day_Glasser}, $\rho_i$ is the
population-weighted average of the densities of the counties of location $i$,
and $\rho$ the population-weighted average of the $\rho_i$; county populations
on July 1, 2020 (\fref{f-co}), land areas (\fref{f-gaz}). Then $\rho=2253$ and
$L_i\in[0.33,1.27]$. The supplement can also be read with $\rho$ the
population-weighted average of the state densities population$/$area, which
gives larger $L_i$. (\fref{f-L})

\subsubsection{Decay with distance}\label{sec-b}
$b_q$ solves
$\sum_lU_{ql}=\frac1m\sum_{r,k}e^{-b_qd_{rk}}$~\cite[Section~2.2.2.2]{2017Feng_Hill_Curns_Glasser};
computed from $U$ and $d$, $b_q\in[0.0009,0.0087]$ per mile.

\subsubsection{Contact matrix}\label{sec-C}
The mixing matrix $C=(c_{ij,pq})$, of order $mn=816$: $c_{ij,pq}$ is the fraction of
the contacts of a person of location $i$ and age group $j$ that are with persons of
location $p$ and age group $q$, so every row of $C$ adds up to $1$. The number of
contacts is given separately, by the contact rates $a_{ij}=L_i\sum_lU_{lj}$.

\emph{Source.} The mixing matrix is defined in~\cite[Section~2.4, pp.~5--6]{2023Glasser_Feng}.
For age groups, $c_{jq}=U_{jq}/a_j$ with $a_j=\sum_kU_{jk}$, so that every row adds up to
$1$, the condition $\sum_qc_{jq}=1$ on mixing matrices stated
in~\cite[Section~2.1]{2017Feng_Hill_Curns_Glasser}. For locations alone, the contacts are
$C_{ip}=a_i\,e^{-bd_{ip}}/\sum_ke^{-bd_{ik}}$ and the mixing is
$c_{ip}=C_{ip}/\sum_kC_{ik}$. For locations and age groups,
\cite[Section~2.4, p.~6]{2023Glasser_Feng} gives
\begin{equation}\label{contacts-glasser}
c_{l_ia_j,\,l_pa_q}=\frac{c^{(p)}_{a_ja_q}\,e^{-b_{a_q}d_{l_il_p}}}
{\sum_{r=1}^{n}\sum_{s=1}^{m}c^{(s)}_{a_ja_r}\,e^{-b_{a_r}d_{l_il_s}}},
\qquad
c^{(p)}_{a_ja_q}=c_{a_ja_q}\,\frac{\ln\left(\rho^{(p)}\right)}{\ln\left(\rho\right)}:
\end{equation}
the age-specific mixing within location $p$ is multiplied by the ratio of the natural
logarithms of the location and national population densities, our $L_p$
(Section~\ref{sec-L}), and between locations the contacts decline exponentially with
distance at the rate of the age group met, computed as
in~\cite{2017Feng_Hill_Curns_Glasser}, our $b_q$ (Section~\ref{sec-b}).
By~\cite[Supplement, Section~2a]{2026Feng_Day_Glasser}, the mixing of the model
of~\cite{2026Feng_Day_Glasser} is the mixing used in~\cite{2023Glasser_Feng}, except that
the densities $\rho_i$ are computed from the county densities (Section~\ref{sec-L}).

\emph{The formula used.} In our notation, with $l_i\to i$, $a_j\to j$,
\eqref{contacts-glasser} is
\begin{equation}\label{contacts}
c_{ij,pq}=\frac{L_pU_{jq}e^{-b_qd_{ip}}}{\sum_{p',q'}L_{p'}U_{jq'}e^{-b_{q'}d_{ip'}}}:
\end{equation}
we use the contacts $U_{jq}$ in place of $c_{jq}=U_{jq}/a_j$; the factor $1/a_j$ depends
only on the row $(i,j)$ and cancels between the numerator and the denominator, so
\eqref{contacts} and~\eqref{contacts-glasser} give the same matrix. A person of $(i,j)$
meets age group $q$ in proportion to $U_{jq}$, location $p$ in proportion to its activity
$L_p$, and the contacts decline with the distance $d_{ip}$ at the rate $b_q$ of the age
group met. For fixed $(i,j)$ the denominator is the sum of the numerator over all $(p,q)$,
so every row adds up to $1$. For $n=1$ and $L_p=1$, \eqref{contacts} is the $c_{ip}$
given for locations alone.

\emph{The file.} $C$ is formed from $U$, $L$, $d$ and $b$ alone by
\src{build_C_51x16.py} (\fref{f-C}), and \src{seir7_ode.py} loads it.
Its entries are in $[1.5\cdot10^{-23},\,0.34]$, every row adds up to $1$ (to
$3\cdot10^{-15}$), and $C\bm x$ agrees with its computation without forming $C$
(Appendix~B of \src{stratified_seir_ngm.tex}) to $2\cdot10^{-15}$.
The rank and the nonzero entries of $C$, and of $C$ after dropping its entries below a
threshold ($\epsilon=2^{-52}$; the rank is the number of singular values, computed by the
SVD, larger than $816\,\epsilon\,\sigma_{\max}$):
\begin{center}
\small
\begin{tabular}{lcc}
 & rank & nonzeros\\\hline
$C$ & 816 & full\\
drop $|c_{ij,pq}|<16\epsilon$ & 816 & $99.62\%$\\
drop $|c_{ij,pq}|<\sqrt{\epsilon}$ & 816 & $95.86\%$
\end{tabular}
\end{center}

\subsubsection{Stringency}\label{sec-s}
$s_i(t)=1-\mathrm{SI}_i(t)/100$, the complement of the
Oxford stringency index of location $i$~\cite[Section~4.1]{2023Glasser_Feng}:
the daily \src{StringencyIndex_Average} of the US states (\src{STATE_TOTAL})
in~\fref{f-ox}, 1 January 2020 -- 31 December 2022, linear between days and
constant after the last. At $t=251$, $s_i\in[0.20,0.62]$. (\fref{f-s})

\subsection{Age-specific values}\label{sec-age}

From Table~S2 of~\cite{2026Feng_Day_Glasser}, which is Table~2
of~\cite{2023Glasser_Feng}: the proportion with mild if any symptoms $p_j$ and
the proportion hospitalized if infected $d_j$~\cite{2020Clark_Jit}, the infection
fatality ratio $\mathrm{IFR}_j$~\cite{2020Levin_Hanage}, and the compliance $c_j^0$,
estimated in~\cite[Section~4.1]{2023Glasser_Feng} from~\cite{2020Jones_YouGov}
(\fref{f-t2}). From these, $q_j=1-(1-p_j)/2$ \cite[Table~2]{2026Feng_Day_Glasser};
$\delta_{s,j}=\gamma_{s\to h}d_j$ and $\delta_{m,j}=\delta_{s,j}/2$ (Table~S1);
$\mathrm{SFR}_j=\mathrm{IFR}_j(1-p_j)$ (Supplement, Section~2); and
$\theta_{h,j}=1-0.7\,\mathrm{SFR}_j$, $\theta_{s,j}=1-0.3\,\mathrm{SFR}_j$, from the
$70\%$ of deaths among inpatients in~\cite[Section~3]{2023Glasser_Feng}.

The probabilities of infection on contact $\beta_j^0$ are the values for 26 November
2020 shown in~\cite[Fig.~4]{2023Glasser_Feng}, used there for transmission
modeling~\cite[Section~2.3]{2023Glasser_Feng}. The figure gives no table; we
digitized the curve at the 16 age groups (axis calibrated on the tick marks,
values read at the centers of the markers; resolution $5.7\cdot10^{-4}$, accuracy
about $\pm0.003$) (\fref{f-beta}):
\begin{center}
\small
\begin{tabular}{l*{8}{c}}
age & 0--4 & 5--9 & 10--14 & 15--19 & 20--24 & 25--29 & 30--34 & 35--39\\
$\beta_j^0$ & 0.121 & 0.078 & 0.064 & 0.054 & 0.080 & 0.071 & 0.073 & 0.067\\\hline
age & 40--44 & 45--49 & 50--54 & 55--59 & 60--64 & 65--69 & 70--74 & 75+\\
$\beta_j^0$ & 0.067 & 0.078 & 0.078 & 0.095 & 0.142 & 0.170 & 0.187 & 0.333
\end{tabular}
\end{center}

\subsection{Constants}\label{sec-const}

From the supplement of~\cite{2026Feng_Day_Glasser}, per day:
$\xi=\gamma_{s\to h}=1/2$, $\gamma_a=1/7$, $\gamma_s=\gamma_m=1/5$, $\gamma_h=1/10$,
$\omega=\omega_v=1/365$, $b_S=b_I=0.31$ (Table~S1);
$\eta_a,\eta_{ps},\eta_{pm},\eta_s,\eta_m,\eta_h=0.5,1.25,1,1,0.75,0.1$ (Section~2).

\subsection{Variants}\label{sec-var}

The values that change with the dominant variant are set as in Table~S4
of~\cite{2026Feng_Day_Glasser} (a blank there keeps the previous value; $k$ for
Wuhan from Table~S4, where Table~S1 gives $1/4$):
\begin{center}
\small
\begin{tabular}{lccccc}
 & Wuhan & Alpha & Delta & BS1 & BA5\\\hline
factor of $\beta_j^0$ & 1 & 1.25 & 1.5 & 2 & 1\\
$\varepsilon_1=\varepsilon_2=\varepsilon_3$ & 0.9 & 0.8 & 0.7 & 0.6 & 0.5\\
factor of $c_j^0$ & 1 & 0.8 & 0.6 & 0.4 & 0.2\\
$1/k$ (days) & 5 & 5 & 4 & 4 & 3\\
factor of $d_j$ and $1-\theta_{w,j}$ & 1 & 1 & 1 & 0.67 & 0.67\\\hline
from (date) & & 3 Apr 2021 & 3 Jul 2021 & 25 Dec 2021 & 2 Jul 2022\\
from ($t$) & & 459 & 550 & 725 & 914
\end{tabular}
\end{center}
A variant is taken to apply from the end of the first week in which its share of
the US sequences exceeds $1/2$, in the weighted estimates of~\fref{f-var} (for
each week the estimates of the latest creation date; lineages, named as in the CDC data,
grouped as Alpha: B.1.1.7, Q; Delta: B.1.617.2, AY; BS1: B.1.1.529, BA.1; BA5: BA.5, BQ,
BF, BE). The names Wuhan, Alpha, Delta, BS1, BA5 are the columns of Table~S4, which
gives no lineages; the lineages of each name and the threshold $1/2$ are our choice.
(\fref{f-vd})

\subsection{Data}\label{sec-inputs}

\subsubsection{Initial state}\label{sec-x0}
% The state at $t=251$. The data: the seroprevalence of the
% CDC Nationwide Commercial Laboratory Seroprevalence
% Survey~\cite{2021Bajema_Wiegand,2023CDC_sero} (\fref{f-sero}), rounds 3
% (26 August -- 9 September 2020) and 4 (9 -- 18 September 2020), by location and
% the 4 age groups 0--17, 18--49, 50--64, 65+. The rules, in place of the
% procedure of~\cite[Section~3]{2023Glasser_Feng}: (1) from 4 to 16 age groups,
% each 5-year group takes the value of the survey group that contains it, and
% 15--19 takes $3/5$ of 0--17 and $2/5$ of 18--49; a missing or suppressed value is
% replaced by the national one of the same round and age group; (2) the
% seropositive fraction $r_{ij}$ is that of round 3, and
% $\lambda_{ij}=\max\left\{-\ln\left((1-r^{(4)}_{ij})/(1-r^{(3)}_{ij})\right)/11.5,\ 10^{-4}\right\}$
% per day; (3) $R=rN$, and $E_s=\lambda S/k$, $I_a=kpE_s/\gamma_a$,
% $I_{ps}=k(1-p)E_s/\xi$, $I_s=\xi I_{ps}/(\gamma_s+\delta_s)$ (the quasi-steady
% form of the 12-day sums of~\cite[Section~3]{2023Glasser_Feng}), with $S$ such
% that the 13 compartments add up to $N_{ij}$, and the other compartments $0$.
% This gives $6.6\%$ seropositive. 

We use what is in (\fref{f-x0}).

\subsubsection{Weekly doses}\label{sec-doses}
% Made up, by location, age group and dose, from the start of vaccination
% (\fref{f-doses}). The rules below are applied over 126 weeks, through the end of the Public
% Health Emergency, 11 May 2023~\cite[Section~2.2, p.~4]{2026Feng_Day_Glasser}, and written to
% three files with the same values in their common weeks:
% \src{mock/doses_51x16x16_mock_126wk.xlsx}, all 126 weeks;
% \src{mock/doses_51x16x16_mock_55wk.xlsx}, the first 55, as many as the reported doses (the
% default); and \src{mock/doses_51x16x16_mock_15wk.xlsx}, the first 15, for a short test run.
% Age group $j$ receives $\mathrm{cov}_jN_{ij}$ first doses evenly over 12 weeks from its
% start week ($\mathrm{cov}=0.9$ from week 1 for 65+; $0.8$ from week 9 for 50--64;
% $0.7$ from week 15 for 20--49; $0.55$, $0.45$, $0.3$, $0.1$ from weeks 17, 22, 47,
% 79 for 15--19, 10--14, 5--9, 0--4), and after these 12 weeks $0.002\,N_{ij}$ first
% doses per week until the end, so that every location has first doses in every
% week; second doses $0.9$ times the first doses four
% weeks earlier; third doses half of all second doses of the group, evenly over 20
% weeks from week 40 (65+) or 48 (20--64); each location, dose and week multiplied
% by $e^{0.1z}$, $z$ standard normal; then each $D_{k,ij}$ rounded to the nearest
% integer. To be replaced by the doses reported to the
% CDC~\cite[Section~2.2]{2026Feng_Day_Glasser}.

We use what is in (\fref{f-doses}).

The file~\fref{f-doses} is organized as a table with three header rows: the first,
``Week 1'', \dots, ``Week $W$'', each in the middle one of its three columns (week $w$
starts on day $t=t_{\mathrm{vacc}}+7(w-1)$, Section~\ref{sec-params-const}); in the second, ``Dose 1'', ``Dose 2'', ``Dose 3+'' under each
week. The column ``State'' and one row for each of the
51 locations, the doses $D_{k,i}=\sum_jD_{k,ij}$ of all ages. 
An empty cell below the two header rows and right of the key columns is read as $0$, and each one is reported in the log with
its cell, location, week and dose (\src{read_sheet} in \src{doses_51x16x16.py});

\section{The model and the computation}\label{sec-model-comp}

The program and its files, in \src{MathBio/code/python/}. Its inputs are the built
files \fref{f-N}--\fref{f-doses} (Section~\ref{sec-built}).
\begin{enumerate}[label=F2-\arabic*, ref=F2-\arabic*, series=program, leftmargin=*]
\item\label{f-prog} \src{seir7_ode.py}: the program, run with the arguments
\texttt{simulate optimal} (Section~\ref{sec-comp}) or \texttt{simulate eigenvector}
(Section~\ref{sec-ev}).
\item\label{f-alt} \src{alternating.py}: the optimization of each week, called
by~\fref{f-prog}.
\end{enumerate}
The computation is run, and after an interruption resumed, from
\src{MathBio/weekly_simulation/} by
\begin{center}
\small
\texttt{./run.sh DOSES INITIAL}\qquad\texttt{RESUME=1 ./run.sh DOSES INITIAL}
\end{center}
with the dose file and the initial state, and the parameter set by \texttt{PARAMS=NAME}
(default \texttt{2021}); \src{run.sh} runs \texttt{simulate eigenvector}, then
\texttt{simulate optimal}, the figures and \src{weekly_report.py}. The results are in
\src{weekly_simulation/results/<tag>/}, with \src{<tag>} the names of the initial state and
dose files without extension and the parameter set, joined by \src{__}:
\begin{enumerate}[label=F2-\arabic*, ref=F2-\arabic*, resume=program, leftmargin=*]
\item\label{f-out-o} \src{simulate_weekly_51x16_optimal__<tag>.csv}: the national
totals of the 13 compartments at the start of each week.
\item\label{f-out-rho} \src{simulate_weekly_51x16_optimal__<tag>_rho.csv}: the
spectral radius of each week at the solution (and, in columns \texttt{rho\_uniform},
\texttt{rho\_eigenvector} and \texttt{rho\_eigenvector\_run}, at the three starts of the
minimization: the uniform allocation, the rates of the eigenvector strategy, and the rates
of the eigenvector run in the same week, divided among the groups in proportion to
$\sigma_{k,p}\chi_{k,p}$).
\item\label{f-chi-ev} \src{simulate_weekly_51x16_eigenvector__<tag>_chi.npz}: the rates of
the eigenvector run in every week, the third start of the optimal run.
\item\label{f-ckpt} \src{simulate_optimal__<tag>_checkpoint.npz}: after every
week, the state, the week and the results so far; with \texttt{RESUME=1} the computation
continues from the last completed week.
\end{enumerate}
The inputs \fref{f-N}--\fref{f-doses} are rebuilt from the downloads by the scripts
named in Section~\ref{sec-built}.

\subsection{The model}\label{sec-model}

The system of~\cite[Section~2.1 and Fig.~1]{2026Feng_Day_Glasser}, with the
location $i$ and the age group $j$ written separately. Each group $(i,j)$ has
13 compartments: susceptible $S$; vaccinated once $V$; partially immune $P$;
exposed $E_s$ (naive) and $E_p$ (partially immune); mild $I_a$; pre-severe and
severe $I_{ps}$, $I_s$; pre-moderate and moderate $I_{pm}$, $I_m$; hospitalized
$I_h$; recovered $R$; dead $M$. The force of infection on group $(i,j)$ is
\begin{equation}\label{lambda}
\lambda_{ij}=\psi_{S,j}\,s_i\,a_{ij}\,\beta_j\sum_{p,q}c_{ij,pq}\,\psi_{I,q}\,
\frac{\sum_x\eta_xI_{x,pq}}{N_{pq}},\qquad x\in\{a,ps,pm,s,m,h\},
\end{equation}
with $a_{ij}$, $C=(c_{ij,pq})$ and $s_i(t)$ as in Section~\ref{sec-loc}, and $\psi_{S,j}=1-c_jb_S$, $\psi_{I,j}=1-c_jb_I$,
$c_j=c_j^0$ and $\beta_j=\beta_j^0$ times the factors of the variant at $t$
(Sections~\ref{sec-age} and~\ref{sec-var}). For every group (indices $i,j$
omitted; $\lambda=\lambda_{ij}(t)$, $\varepsilon=\varepsilon(t)$ of the variant,
$\chi_k=\chi_{k,ij}$ the vaccination rates of dose $k$),
\begin{equation}\label{ode}
\begin{aligned}
S'&=\omega P+\omega_vV-(\lambda+\chi_1)S, &
V'&=\chi_1S-\left((1-\varepsilon)\lambda+\omega_v+\varepsilon\chi_2\right)V,\\
P'&=\omega R-(\omega+\varepsilon\chi_3+\lambda)P, &
E_s'&=\lambda S-kE_s,\\
E_p'&=(1-\varepsilon)\lambda V+\lambda P-kE_p, &
I_a'&=kp_jE_s+kq_jE_p-\gamma_aI_a,\\
I_{ps}'&=k(1-p_j)E_s-\xi I_{ps}, &
I_{pm}'&=k(1-q_j)E_p-\xi I_{pm},\\
I_s'&=\xi I_{ps}-(\gamma_s+\delta_{s,j})I_s, &
I_m'&=\xi I_{pm}-(\gamma_m+\delta_{m,j})I_m,\\
I_h'&=\delta_{s,j}I_s+\delta_{m,j}I_m-\gamma_hI_h, &
M'&=(1-\theta_{s,j})\gamma_sI_s+(1-\theta_{h,j})\gamma_hI_h,
\end{aligned}
\end{equation}
\begin{equation*}
R'=\gamma_aI_a+\theta_{s,j}\gamma_sI_s+\gamma_mI_m+\theta_{h,j}\gamma_hI_h
+\varepsilon\chi_2V+\varepsilon\chi_3P-\omega R,
\end{equation*}
$13mn=10\,608$ unknowns, coupled only through $\lambda$. The 13 right hand sides
of a group add up to $0$, so each group keeps its size $N_{ij}$.

\subsection{The computation}\label{sec-comp}

\begin{enumerate}[leftmargin=*]
\item Solve~\eqref{ode} from $t=\RunTzero$, with the initial state~\fref{f-x0}, to
$t=\RunTvacc$, the start of vaccination of the parameter set \texttt{\RunParams}, with $\chi=0$.
\item For each week $[\tau,\tau+7]$, $\tau=\RunTvacc,\dots,\RunTlast$ (\RunWeeks{} weeks, to
$t=\RunTend$, \RunTendDate), with
$S$, $V$, $P$ the state at $\tau$ and $D_{k,i}$ the doses of location $i$ in the
week in~\fref{f-doses} (Section~\ref{sec-doses}):
\begin{enumerate}[label=(\alph*)]
\item compute the vaccination rates $\chi_{k,ij}$ for the week (the same on every
day of the week; they change only from one week to the next): minimize the
spectral radius of the next generation matrix over $\chi$, subject to
$\sum_j\sigma_{k,ij}\chi_{k,ij}=\frac17D_{k,i}$ for every location $i$ and
dose $k$ (the doses of the location, all ages, at most $0.99\sum_j\sigma_{k,ij}/7$),
with $\sigma_1=S$, $\sigma_2=V$, $\sigma_3=P$, only the groups with
$\sigma_{k,ij}\ge1$ (at least one eligible person) receiving dose $k$, and the bound
$0\le\chi_{k,ij}\le\chi_{\max}=-\frac17\ln(0.01)\approx0.658$ per day (at most $99\%$
of a group in a week) as a constraint of the minimization
(\src{stratified_seir_ngm.tex}, problem (1.2), Section~6 and its Section~6.1);
\item solve~\eqref{ode} on $[\tau,\tau+7]$ with these rates, which gives the state
at $\tau+7$;
\item advance to the next week, $\tau\leftarrow\tau+7$.
\end{enumerate}
\end{enumerate}
The system is solved by the Runge--Kutta method of order 8 of Dormand and
Prince~\cite{1980Dormand_Prince} (\texttt{DOP853} of
\texttt{scipy.integrate.solve\_ivp}~\cite{2020Virtanen_scipy}), restarted at the
days at which the variant changes.

\begin{remark}[The actual strategy]\label{rem-actual}
In~\cite{2026Feng_Day_Glasser} the optimal allocation is compared with the
actual one, the doses reported to the CDC by location, age group and dose. In
the actual strategy no optimization is done: in each week, the fraction
$P_{k,ij}=\min\{D_{k,ij}/X_{k,ij}(\tau),0.99\}$ of the eligible people
$X_1=S$, $X_2=V$, $X_3=P$ of group $(i,j)$ is vaccinated, that is,
$\chi_{k,ij}=-\frac17\ln\left(1-P_{k,ij}\right)$, and the system~\eqref{ode} is
solved with these rates. With the doses by age group
of~\cite[Section~2.2]{2026Feng_Day_Glasser}, it gives the reference against which
the optimal allocation is measured (illnesses, hospitalizations and deaths
averted). It is not computed here. 
\end{remark}

\subsection{The eigenvector strategy}\label{sec-ev}

The eigenvector strategy is the alternative strategy
of~\cite[Section~2.2 and Appendix~A]{2026Feng_Day_Glasser}: each week the doses of each
location and dose, summed over the ages, go to its eligible age groups in proportion to the
left eigenvector of the dominant eigenvalue of the next generation matrix. It is compared
with the optimal strategy, each with its own solution of~\eqref{ode} from
$t=t_{\mathrm{vacc}}$ and the same doses $b_{k,i}$ per day of each location, dose and week
(Section~\ref{sec-comp}).

\paragraph{The rates.} With the notation of \src{stratified_seir_ngm.tex}, Sections~1--2,
let $\bm u>0$ be the left Perron vector of $H_1\operatorname{diag}(\bm g)$ and
$J_{k,i}=\{p\in I_i:\ \sigma_{k,p}\ge1\}$ the eligible groups of location $i$ for dose
$k$. The weights are
\begin{equation}\label{w-ev}
w_{1,p}=\frac{u_p}{g_p},\qquad w_{2,p}=w_{3,p}=r_p\,\frac{u_p}{g_p},
\end{equation}
and the rates of the week are
\begin{equation}\label{chi-ev}
\chi_{k,p}=\frac{b_{k,i}}{\sigma_{k,p}}\,\frac{w_{k,p}}{\sum_{l\in J_{k,i}}w_{k,l}},\quad
p\in J_{k,i},\qquad\chi_{k,p}=0,\quad p\in I_i\setminus J_{k,i},
\end{equation}
followed by the bound $\chi_{k,p}\le\chi_{\max}$ as in (C4) below. By Lemma~\ref{lem-ev},
$\rho\,w_{1,p}$ and $\rho\,w_{2,p}$ are the entries $2p-1$ and $2p$ of the left Perron vector
of the next generation matrix $H$, the matrix defined in~\cite[Appendix~A]{2026Feng_Day_Glasser},
whose rows $2p-1$ and $2p$ are the infections of group $p$ that start in $S$ (through $E_s$)
and in $V$ or $P$ (through $E_p$).

\paragraph{Our choices.} \cite{2026Feng_Day_Glasser} does not specify the following; we take:
\begin{enumerate}[label=(C\arabic*), leftmargin=*]
\item The time: $\bm u$ and $\bm g$ at $\tau$, with the parameters and the state at $\tau$
and the rates $\chi_{k,p}$ of the previous week ($\chi=0$ in the first week).
\item The entries: entry $2p-1$ for the first dose, given to $S$; entry $2p$ for the second
and the third dose, given to $V$ and $P$.
\item Doses and rates: the doses $\sigma_{k,p}\chi_{k,p}$, not the rates, are in proportion
to the weights, and they add up to $b_{k,i}$, the constraint of the optimal problem; so the
strategies give the same doses per day. The conversion $\chi=-\frac17\ln(1-\nu)$ of
the actual strategy (Remark~\ref{rem-actual}) is not used.
\item The bound: a rate above $\chi_{\max}$ is lowered to $\chi_{\max}$, and the doses above
it go to the other eligible groups of the location below it, in proportion to their doses
(\src{stratified_seir_ngm.tex}, Lemma~6.5: this ends with all rates $\le\chi_{\max}$ and
the same doses).
\end{enumerate}

\begin{lemma}\label{lem-ev}
Let $C>0$, $\bm d_c,\bm\psi_I,\bm q_s,\bm g>0$, $r_p=q_{m,p}/q_{s,p}$, and let $\bm u>0$ be
the left Perron vector of $H_1\operatorname{diag}(\bm g)$, with eigenvalue $\rho$. With
$H=X_1X_2$, $X_1=\operatorname{diag}(\bm d)\mathcal{I}$,
$X_2=\operatorname{diag}(\bm d_c)C\mathcal{I}^T\Psi_IQ$ (\src{stratified_seir_ngm.tex},
Theorem~2.1), $\bm w^T=\bm u^TX_2$ is a left Perron vector of $H$, and
$w_{2p-1}=\rho\,u_p/g_p$, $w_{2p}=r_pw_{2p-1}$, $p=1:mn$.
\end{lemma}
\begin{proof}
$H_1=\operatorname{diag}(\bm d_c)C\operatorname{diag}(\bm\psi_I\odot\bm q_s)>0$, so
$H_1\operatorname{diag}(\bm g)>0$ is irreducible and, by the Perron--Frobenius theorem,
$\rho>0$ is a simple eigenvalue with a left eigenvector $\bm u>0$. Since
$X_2X_1=H_1\operatorname{diag}(\bm g)$, $\bm w^TH=\bm u^TX_2X_1X_2=\rho\bm u^TX_2
=\rho\bm w^T$, and $\bm w\ne0$ since $\bm w^TX_1=\rho\bm u^T\ne0$. Entry $2p-1$ of
$\bm u^T\operatorname{diag}(\bm d_c)C\mathcal{I}^T\Psi_IQ$ is
$\left(C^T\operatorname{diag}(\bm d_c)\bm u\right)_p\psi_{I,p}q_{s,p}=(H_1^T\bm u)_p$, and entry
$2p$ is the same with $q_{m,p}=r_pq_{s,p}$ in place of $q_{s,p}$. From
$\bm u^TH_1\operatorname{diag}(\bm g)=\rho\bm u^T$, $(H_1^T\bm u)_pg_p=\rho u_p$. So $\bm w>0$,
and it is the left Perron vector of $H\ge0$: a nonnegative matrix has, up
to a factor, no other positive left eigenvector.
\end{proof}

For each strategy and week, $\rho$ is computed at $\tau$ for the rates applied in the week,
with the state of that strategy. The strategies are compared by $\rho$, by the infectious,
hospitalized and dead over time, and by the doses given to each age group.

\section[The two parameter sets: Unknown 2020/21 and Feng, Day, Glasser 2026; Glasser et al.\ 2023]{The two parameter sets: Unknown 2020/21 and~\cite{2026Feng_Day_Glasser},~\cite{2023Glasser_Feng}}\label{sec-params}

The model is computed with one of two parameter sets: the values of Unknown 2020/21
(set \src{2021}, the default) and the values used in~\cite{2026Feng_Day_Glasser}
and~\cite{2023Glasser_Feng}, as described in Section~\ref{sec-data} (set \src{feng2026}). Both are
in \src{weekly_simulation/input/parameters.xlsx}, one column per set (sheets \src{scalars},
\src{by_age}, \src{over_time}); the set is chosen by \src{params=2021} or
\src{params=feng2026} in \src{seir7_ode.py}, \src{ngm_seir7.py} and
\src{plot_weekly_optimal.py}, and its name ends the names of the output files.

\subsection{Constants of the two sets}\label{sec-params-const}

Per day, unless stated otherwise:
\begin{center}
\small
\begin{tabular}{lcc}
 & Unknown 2020/21 & \cite{2026Feng_Day_Glasser}, \cite{2023Glasser_Feng}\\\hline
$\xi$, $\gamma_{s\to h}$ & $1/2$, $1/2$ & $1/2$, $1/2$\\
$\gamma_a$, $\gamma_s$, $\gamma_m$, $\gamma_h$ & $1/7$, $1/5$, $1/5$, $1/10$ & $1/7$, $1/5$, $1/5$, $1/10$\\
$\omega=\omega_v$ & $1/182.5$ & $1/365$\\
$b_S=b_I$ & 0.31 & 0.31\\
$\eta_a,\eta_{ps},\eta_{pm},\eta_s,\eta_m,\eta_h$ & 0.5, 1.25, 1, 1, 0.75, 0.1 & 0.5, 1.25, 1, 1, 0.75, 0.1\\
deaths of $\mathrm{SFR}_j$: outside, in hospital & 0.33, 0.67 & 0.3, 0.7\\
initial state, $t_0$ & $t=347$ (12 Dec 2020) & $t=251$ (7 Sep 2020)\\
start of vaccination & $t=354$ (19 Dec 2020) & $t=349$ (14 Dec 2020)
\end{tabular}
\end{center}
That is, $\theta_{s,j}=1-0.33\,\mathrm{SFR}_j$ and $\theta_{h,j}=1-0.67\,\mathrm{SFR}_j$ in the
set of Unknown 2020/21. In both sets $q_j=1-(1-p_j)/2$, $\delta_{s,j}=\gamma_{s\to h}d_j$ and
$\delta_{m,j}=\delta_{s,j}/2$. The doses of week $w$ are given from
$t=t_{\mathrm{vacc}}+7(w-1)$, with $t_{\mathrm{vacc}}$ the start of vaccination of the
set.

\subsection{Age-specific values of the two sets}\label{sec-params-age}

The values $p_j$, $d_j$ and $\mathrm{IFR}_j$ are the same in both sets
(Section~\ref{sec-age}); the compliance $c_j^0$ of Unknown 2020/21 is given with 16 digits and
agrees with the compliance described in Section~\ref{sec-age} to the 9 digits given there.
The probabilities of infection on contact $\beta_j^0$ differ (the values from
Fig.~4 of~\cite{2023Glasser_Feng} are described in Section~\ref{sec-age}):
\begin{center}
\small
\begin{tabular}{l*{8}{c}}
age & 0--4 & 5--9 & 10--14 & 15--19 & 20--24 & 25--29 & 30--34 & 35--39\\\hline
Unknown 2020/21 & 0.1241 & 0.0815 & 0.0645 & 0.0511 & 0.0772 & 0.0762 & 0.0733 & 0.0683\\
from Fig.~4~\cite{2023Glasser_Feng} & 0.1211 & 0.0779 & 0.0637 & 0.0535 & 0.0796 & 0.0711 & 0.0728 & 0.0671\\\hline
age & 40--44 & 45--49 & 50--54 & 55--59 & 60--64 & 65--69 & 70--74 & 75+\\\hline
Unknown 2020/21 & 0.0689 & 0.0794 & 0.0793 & 0.0951 & 0.1426 & 0.1708 & 0.1881 & 0.3305\\
from Fig.~4~\cite{2023Glasser_Feng} & 0.0671 & 0.0779 & 0.0779 & 0.0950 & 0.1422 & 0.1701 & 0.1871 & 0.3327
\end{tabular}
\end{center}
The values from Fig.~4~\cite{2023Glasser_Feng} differ from the values of Unknown 2020/21 by $-6.7\%$ (25--29) to $+4.7\%$
(15--19), and by less than $1\%$ from age 55.

\subsection{Values that change over time in the two sets}\label{sec-params-var}

In the set of Unknown 2020/21 each value has its own days of change, and the values after
1 July 2021 are kept to the end (there are no values for BS1 and BA5); a value given
from day $t$ holds for the days after $t$:
\begin{center}
\small
\begin{tabular}{lcccc}
 & start & from $t=398$ & from $t=460$ & from $t=548$\\
 & & (1 Feb 2021) & (4 Apr 2021) & (1 Jul 2021)\\\hline
factor of $\beta_j^0$ & 1 & 1.25 & & 1.5\\
$\varepsilon_1=\varepsilon_2=\varepsilon_3$ & 0.9 & 0.8 & & 0.7\\
factor of $c_j^0$ & 1 & & 0.8 & \\
$1/k$ (days) & 3 & & & \\
factor of $d_j$ and $1-\theta_{w,j}$ & 1 & & &
\end{tabular}
\end{center}
A blank keeps the previous value. In the set used in~\cite{2026Feng_Day_Glasser}, as
described in Section~\ref{sec-var}, all values change together on the days of the variants,
$t=459$, $550$, $725$, $914$. The two sets differ in the day of the first change ($t=398$
for $\beta$ and $\varepsilon$, $t=460$ for the compliance, against $t=459$ for all), in
the factor of $c_j^0$ under Delta (0.8 against 0.6), in $k$ ($1/3$ throughout, against
$1/5$, $1/4$ from Delta and $1/3$ from BA5), and after Delta.

\bibliographystyle{plainurl}
\bibliography{bib_mb}
\end{document}
